% !TEX program = lualatex
\documentclass[professionalfonts]{beamer}
\newcommand{\slidemode}{1}  % カラー投影モード
\usepackage{beamer_template}
\usepackage{enumerate}

\newcommand{\LectureNum}{6}
\newcommand{\LectureTitle}{たたみ込み}
\newcommand{\CourseName}{応用数学}
\renewcommand{\TermName}{前期}

\begin{document}

\begin{frame}{第6回　たたみ込み}
  \begin{block}{今回の内容}
  たたみ込み
  \end{block}
  \vfill\hfill{\scriptsize \textcolor{gray}{第2章 ラプラス変換　§2}}
\end{frame}

\begin{frame}{定義：たたみこみ（合成積）}
  \begin{block}{}
  \[
    (f * g)(t) = \int_{0}^t f(\tau)g(t-\tau) d\tau
  \]
  {\footnotesize \textcolor{gray}{$f(t), g(t)$ はいずれも $t \geq 0$ で定義された関数}}
  \end{block}
\end{frame}

\begin{frame}{定理：交換法則}
  \begin{block}{}
  \[
    f(t) * g(t) = g(t) * f(t)
  \]
  \end{block}
\end{frame}

\begin{frame}{定理：分配法則}
  \begin{block}{}
  \[
    f * (g_{1} + g_{2}) = f*g_{1} + f*g_{2}
  \]
  \end{block}
\end{frame}

\begin{frame}{定理：たたみこみのラプラス変換（たたみこみ定理）}
  \begin{block}{}
  \[
    L[f(t) * g(t)] = L[f(t)] \cdot L[g(t)] = F(s) \cdot G(s)
  \]
  \[
    L^{-1}[F(s) \cdot G(s)] = f(t) * g(t) = \int_{0}^t f(\tau)g(t-\tau) d\tau
  \]
  \end{block}
\end{frame}

\begin{frame}{例題5}
  \begin{exampleblock}{問題}
    $\sin t * \cos t$ を求めよ\\[2pt]
  \end{exampleblock}
\end{frame}

\begin{frame}{補足：積和公式}
  \begin{block}{公式}
  \[
    \sin A \cos B = \frac{1}{2}[\sin(A+B) + \sin(A-B)]
  \]
  \end{block}
  \textbf{導出}：加法定理より
  \begin{align*}
    \sin(A+B) &= \sin A \cos B + \cos A \sin B \\
    \sin(A-B) &= \sin A \cos B - \cos A \sin B
  \end{align*}
  辺々足すと $\cos A \sin B$ が消えて
  \[
    \sin(A+B) + \sin(A-B) = 2\sin A \cos B
  \]
  \textbf{例題5への適用}：$A=\tau,\; B=t-\tau$ とおくと
  \[
    A+B = t,\quad A-B = 2\tau - t
  \]
  \[
    \therefore\quad \sin\tau\cos(t-\tau) = \tfrac{1}{2}[\sin t + \sin(2\tau-t)]
  \]
\end{frame}

\begin{frame}{例題5　解答}
  定義より $\sin t * \cos t = \int_{0}^t \sin \tau \cdot \cos(t-\tau) d\tau$\\[4pt]
  積和公式 $: \sin \tau\cdot \cos(t-\tau) = (1/2)[\sin(\tau+(t-\tau)) + \sin(\tau-(t-\tau))]$\\[4pt]
  $= (1/2)[\sin t + \sin(2\tau-t)]$\\[4pt]
  $\int_{0}^t (1/2)[\sin t + \sin(2\tau-t)]d\tau$\\[4pt]
  $= (1/2)\left\{ [\sin t \cdot \tau]_{0}^t + [-\cos(2\tau-t)/2]_{0}^t \right\}$\\[4pt]
  $= (1/2)\left\{ t\sin t + (-\cos t/2 - (-\cos(-t)/2)) \right\}$\\[4pt]
  $= (1/2)\left\{ t\sin t + 0 \right\}$（$\cos(-t) = \cos t$ より）
  \\[8pt]\textcolor{red}{\textbf{答}}：$\sin t * \cos t = (t/2) \sin t$
\end{frame}

\begin{frame}{演習問題}
  \begin{exampleblock}{自分で解いてみよう}
  \begin{description}
    \item[問5] 次のたたみこみを求めよ\\
    $t^{2} * t$\\
  \end{description}
  \end{exampleblock}
\end{frame}

\begin{frame}{問5　解答}
  $t^{2} * t = \int_{0}^t \tau^{2}(t-\tau) d\tau = \int_{0}^t (\tau^{2}t - \tau^{3}) d\tau$\\[4pt]
  $= t\cdot[\tau^{3}/3]_{0}^t - [\tau^{4}/4]_{0}^t = t^{4}/3 - t^{4}/4 = t^{4}/12$
  \\[8pt]\textcolor{red}{\textbf{答}}：$t^{2} * t = t^{4}/12$
\end{frame}

\begin{frame}{例題6}
  \begin{exampleblock}{問題}
    $f(t) * g(t) = g(t) * f(t)$ を証明せよ\\[2pt]
  \end{exampleblock}
\end{frame}

\begin{frame}{補足：置換積分の積分限界}
  $u = t - \tau$ と置換するときの変数変化を確認する。\\[6pt]
  \begin{block}{置換のルール}
  \begin{tabular}{ll}
    微分 & $u = t-\tau \;\Rightarrow\; du = -d\tau$ \\[4pt]
    限界 & $\tau = 0 \;\Rightarrow\; u = t$ \\
         & $\tau = t \;\Rightarrow\; u = 0$
  \end{tabular}
  \end{block}
  \vspace{4pt}
  よって積分は
  \[
    \int_{\tau=0}^{\tau=t} (\cdots)\,d\tau
    = \int_{u=t}^{u=0} (\cdots)(-du)
    = \int_{0}^{t} (\cdots)\,du
  \]
  \textbf{ポイント}：$-du$ の符号と限界の逆転が打ち消し合い、\\
  結果として符号が変わらない。
\end{frame}

\begin{frame}{例題6　解答}
  $f(t)*g(t) = \int_{0}^t f(\tau)g(t-\tau)d\tau$\\[4pt]
  $u = t-\tau$ と置換 $: d\tau = -du, \tau=0 \to u=t, \tau=t \to u=0$\\[4pt]
  $= \int_{t}^0 f(t-u)g(u)(-du) = \int_{0}^t g(u)f(t-u)du$\\[4pt]
  $= g(t)*f(t)$\hfill$\square$
\end{frame}

\begin{frame}{演習問題}
  \begin{exampleblock}{自分で解いてみよう}
  \begin{description}
    \item[問6] 分配法則 $f*(g_{1}+g_{2}) = f*g_{1} + f*g_{2}$ を証明せよ
  \end{description}
  \end{exampleblock}
\end{frame}

\begin{frame}{問6　解答}
  定義 $: f*(g_{1}+g_{2}) = \int_{0}^t f(\tau)(g_{1}(t-\tau)+g_{2}(t-\tau))d\tau$\\[4pt]
  積分の線形性より $= \int_{0}^t f(\tau)g_{1}(t-\tau)d\tau + \int_{0}^t f(\tau)g_{2}(t-\tau)d\tau$\\[4pt]
  $= f*g_{1} + f*g_{2}$（証明終）
\end{frame}

\begin{frame}{演習問題}
  \begin{exampleblock}{自分で解いてみよう}
  \begin{description}
    \item[問7] $L[t^{2}*t]$ を求めよ。また、問 5 の結果から直接 $L[t^{2}*t]$ を求めよ
  \end{description}
  \end{exampleblock}
\end{frame}

\begin{frame}{問7　解答}
  【たたみこみ定理を使う方法】\\[4pt]
  $L[t^{2}*t] = L[t^{2}]\cdot L[t] = (2/s^{3})\cdot(1/s^{2}) = 2/s^{5}$\\[8pt]
  【問5の結果から直接求める方法】\\[4pt]
  問5より $t^{2}*t = t^{4}/12$ であるから\\[4pt]
  $L[t^{2}*t] = L[t^{4}/12] = (1/12)\cdot4!/s^{5} = 2/s^{5}$
  \\[8pt]\textcolor{red}{\textbf{答}}：$L[t^{2}*t] = 2/s^{5}$
\end{frame}

\begin{frame}{例題7}
  \begin{exampleblock}{問題}
    $L[f(t)] = F(s)$ のとき、関数 $F(s)/(s-\alpha)$ の逆ラプラス変換を求めよ。ただし $\alpha$ は定数とする。\\[2pt]
  \end{exampleblock}
\end{frame}

\begin{frame}{例題7　解答}
  $L^{-1}[F(s)/(s-\alpha)] = L^{-1}[F(s) \cdot 1/(s-\alpha)]$\\[4pt]
  $L^{-1}[1/(s-\alpha)] = e^{\alpha t}$ より、たたみこみ定理を適用\\[4pt]
  $= L^{-1}[F(s)] * e^{\alpha t} = f(t) * e^{\alpha t}$
  \\[8pt]\textcolor{red}{\textbf{答}}：$L^{-1}[F(s)/(s-\alpha)] = \int_{0}^t e^{\alpha(t-\tau)} f(\tau) d\tau$
\end{frame}

\begin{frame}{演習問題}
  \begin{exampleblock}{自分で解いてみよう}
  \begin{description}
    \item[問8] $L[f(t)] = F(s)$ のとき、次の関数の逆ラプラス変換を求めよ。
    ただし $\alpha, \beta$ は定数で、 (2) では $\alpha \neq \beta$ 、 (3) では $\beta \neq 0$ とする。\\
    $(1) F(s)/(s-\alpha)^{2}$\\
    $(2) F(s)/((s-\alpha)(s-\beta))$\\
    $(3) F(s)/((s-\alpha)^{2}+\beta^{2})$\\
  \end{description}
  \end{exampleblock}
\end{frame}

\begin{frame}{問8　解答（1）}
  $F(s)/(s-\alpha)^{2}$ の逆ラプラス変換を求める\\[8pt]
  $F(s)\cdot\dfrac{1}{(s-\alpha)^{2}}$ と分解する\\[8pt]
  $L^{-1}\!\left[\dfrac{1}{(s-\alpha)^{2}}\right] = te^{\alpha t}$ であるから、たたみこみ定理より\\[8pt]
  $L^{-1}\!\left[\dfrac{F(s)}{(s-\alpha)^{2}}\right] = f(t) * te^{\alpha t}$
  \\[8pt]\textcolor{red}{\textbf{答}}：$\displaystyle\int_{0}^t f(\tau)\cdot(t-\tau)\cdot e^{\alpha(t-\tau)}d\tau$
\end{frame}

\begin{frame}{問8　解答（2）}
  $F(s)/((s-\alpha)(s-\beta))$ の逆ラプラス変換を求める（$\alpha \neq \beta$）\\[8pt]
  まず $\dfrac{1}{(s-\alpha)(s-\beta)}$ を部分分数分解する\\[8pt]
  $\dfrac{1}{(s-\alpha)(s-\beta)} = \dfrac{A}{s-\alpha} + \dfrac{B}{s-\beta}$\\[8pt]
  両辺に $(s-\alpha)(s-\beta)$ を掛けて $: 1 = A(s-\beta) + B(s-\alpha)$\\[4pt]
  $s = \alpha$ を代入 $: 1 = A(\alpha-\beta) \;\Rightarrow\; A = \dfrac{1}{\alpha-\beta}$\\[4pt]
  $s = \beta$ を代入 $: 1 = B(\beta-\alpha) \;\Rightarrow\; B = \dfrac{1}{\beta-\alpha} = -\dfrac{1}{\alpha-\beta}$\\[8pt]
  $\therefore\quad \dfrac{1}{(s-\alpha)(s-\beta)} = \dfrac{1}{\alpha-\beta}\!\left(\dfrac{1}{s-\alpha} - \dfrac{1}{s-\beta}\right)$
\end{frame}

\begin{frame}{問8　解答（2）つづき}
  $L^{-1}\!\left[\dfrac{1}{(s-\alpha)(s-\beta)}\right]
   = \dfrac{1}{\alpha-\beta}\!\left(e^{\alpha t} - e^{\beta t}\right)$\\[8pt]
  したがって、たたみこみ定理より\\[8pt]
  $L^{-1}\!\left[\dfrac{F(s)}{(s-\alpha)(s-\beta)}\right]
   = f(t) * \dfrac{e^{\alpha t} - e^{\beta t}}{\alpha-\beta}$
  \\[8pt]\textcolor{red}{\textbf{答}}：$\displaystyle\frac{1}{\alpha-\beta}\int_{0}^t f(\tau)\!\left(e^{\alpha(t-\tau)}-e^{\beta(t-\tau)}\right)d\tau$
\end{frame}

\begin{frame}{補足：移動法則（第一移動定理）}
  \begin{block}{定理}
  \[
    L\bigl[e^{at}f(t)\bigr] = F(s-a), \quad
    L^{-1}\bigl[F(s-a)\bigr] = e^{at}f(t)
  \]
  \end{block}
  \textbf{使い方}：$s$ を $s-\alpha$ に置き換えると、逆変換に $e^{\alpha t}$ が掛かる。\\[6pt]
  \textbf{問8(3)での適用}：\\[4pt]
  $L^{-1}\!\left[\dfrac{\beta}{s^{2}+\beta^{2}}\right] = \sin(\beta t)$ は既知。\\[8pt]
  $s \to s-\alpha$ に移動すると
  \[
    L^{-1}\!\left[\frac{\beta}{(s-\alpha)^{2}+\beta^{2}}\right] = e^{\alpha t}\sin(\beta t)
  \]
  両辺を $\beta$ で割ると
  \[
    L^{-1}\!\left[\frac{1}{(s-\alpha)^{2}+\beta^{2}}\right] = \frac{1}{\beta}e^{\alpha t}\sin(\beta t)
  \]
\end{frame}

\begin{frame}{問8　解答（3）}
  $F(s)/((s-\alpha)^{2}+\beta^{2})$ の逆ラプラス変換を求める（$\beta \neq 0$）\\[8pt]
  $L^{-1}\!\left[\dfrac{1}{(s-\alpha)^{2}+\beta^{2}}\right]$：移動法則より\\[4pt]
  $L^{-1}\!\left[\dfrac{\beta}{s^{2}+\beta^{2}}\right] = \sin(\beta t)$ を $s \to s-\alpha$ に移動\\[8pt]
  $\therefore\quad L^{-1}\!\left[\dfrac{1}{(s-\alpha)^{2}+\beta^{2}}\right] = \dfrac{1}{\beta}e^{\alpha t}\sin(\beta t)$\\[8pt]
  たたみこみ定理より\\[4pt]
  $L^{-1}\!\left[\dfrac{F(s)}{(s-\alpha)^{2}+\beta^{2}}\right] = f(t) * \dfrac{1}{\beta}e^{\alpha t}\sin(\beta t)$
  \\[8pt]\textcolor{red}{\textbf{答}}：$\displaystyle\frac{1}{\beta}\int_{0}^t f(\tau)e^{\alpha(t-\tau)}\sin(\beta(t-\tau))d\tau$
\end{frame}

\begin{frame}{例題8}
  \begin{exampleblock}{問題}
    $\int_{0}^t x(\tau) \sin(t-\tau) d\tau = t^{2}$ を満たす関数 $x(t)$ を求めよ（積分方程式）\\[2pt]
  \end{exampleblock}
\end{frame}

\begin{frame}{補足：積分方程式でたたみこみを見抜く}
  \begin{block}{たたみこみの形}
  \[
    \int_{0}^{t} f(\tau)\,g(t-\tau)\,d\tau = f(t) * g(t)
  \]
  \end{block}
  \textbf{チェックポイント}：被積分関数が $(\tau\text{の関数})\times(t-\tau\text{の関数})$ の積になっているか。\\[8pt]
  \textbf{例題8の場合}
  \[
    \int_{0}^{t} x(\tau)\sin(t-\tau)\,d\tau
  \]
  \begin{itemize}
    \item $f(\tau) = x(\tau)$（求めたい関数）
    \item $g(t-\tau) = \sin(t-\tau)$
  \end{itemize}
  \[
    \therefore\quad \int_{0}^{t} x(\tau)\sin(t-\tau)\,d\tau = x(t)*\sin t
  \]
  両辺をラプラス変換すると $X(s)\cdot L[\sin t]$ と分離できる。
\end{frame}

\begin{frame}{例題8　解答}
  左辺は $x(t) * \sin t$ のたたみこみ\\[4pt]
  両辺をラプラス変換 $: X(s) \cdot L[\sin t] = L[t^{2}]$\\[4pt]
  $X(s) \cdot 1/(s^{2}+1) = 2/s^{3}$\\[4pt]
  $X(s) = 2(s^{2}+1)/s^{3} = 2/s + 2/s^{3}$\\[4pt]
  逆ラプラス変換 $: x(t) = 2 + t^{2}$
  \\[8pt]\textcolor{red}{\textbf{答}}：$x(t) = t^{2} + 2$
\end{frame}

\begin{frame}{演習問題}
  \begin{exampleblock}{自分で解いてみよう}
  \begin{description}
    \item[問9] 次の積分方程式を満たす関数 $x(t)$ を求めよ。
    $x(t) + \int_{0}^t x(t-\tau)e^\tau d\tau = \cos 3t$\\
  \end{description}
  \end{exampleblock}
\end{frame}

\begin{frame}{問9　解答}
  $\int_{0}^t e^\tau x(t-\tau) d\tau = e^t * x(t) = x(t)*e^t$ （たたみこみ）\\[4pt]
  両辺をラプラス変換 $: X(s) + X(s)\cdot1/(s-1) = s/(s^{2}+9)$\\[4pt]
  $X(s)(1 + 1/(s-1)) = s/(s^{2}+9)$\\[4pt]
  $X(s)\cdot s/(s-1) = s/(s^{2}+9)$\\[4pt]
  $X(s) = (s-1)/(s^{2}+9) = s/(s^{2}+9) - 1/(s^{2}+9)$
  \\[8pt]\textcolor{red}{\textbf{答}}：$x(t) = \cos(3t) - (1/3)\sin(3t)$
\end{frame}

\end{document}
